\documentclass[11pt]{article}
\usepackage[margin=2.1cm]{geometry}
\usepackage{booktabs}
\usepackage{amsmath}
\usepackage{enumitem}
\usepackage{xcolor}
\usepackage[T1]{fontenc}
\usepackage{lmodern}

\setlength{\parindent}{0pt}
\setlength{\parskip}{3.6pt}
\newcommand{\code}[1]{\texttt{#1}}
\newcommand{\hr}{\vspace{2pt}\hrule\vspace{6pt}}

\begin{document}

\begin{center}
{\Large\bfseries EE\,362 --- Computer Organization Laboratory}\\[2pt]
{\large\bfseries Lab 3 --- Multiplexers}\\[4pt]
Fall 2026 \quad$\cdot$\quad Lab TA: Gandhi Torres \quad$\cdot$\quad Sections: Mon / Tue, 3:45--5:45~PM, ECIII 237
\end{center}
\hr

\section*{1. Objective}

Build the component that lets a control signal steer a datapath, and learn the
kind of specification table that makes wide designs possible at all.

\begin{quote}
\emph{A multiplexer is how one control bit steers thirty-two bits of data.}
\end{quote}

\section*{2. Background: why the table shrinks}

A 32-bit 3-to-1 multiplexer has 98 inputs. An exhaustive truth table would need
$2^{98}$ rows --- more rows than there are atoms in your body. You are not going
to write that table, and you do not need to.

\textbf{Type C, a compressed table}, enumerates only the \emph{control} inputs and
names a \emph{source} for the data:

\begin{center}
\begin{tabular}{clc@{\hskip 8pt}||@{\hskip 8pt}ll}
\toprule
\# & \code{Sel1} & \code{Sel0} & $Y$ & note \\
\midrule
0 & 0 & 0 & \textbf{A} & \\
1 & 0 & 1 & \textbf{B} & \\
2 & 1 & 0 & \textbf{C} & \\
3 & 1 & 1 & \textbf{0} & \emph{defined default} --- a decision, not an accident \\
\bottomrule
\end{tabular}
\end{center}

Four rows. Complete specification. Rows say \emph{``Y = B''}, naming a source, not
a value. This is not a shortcut --- it is the reason datapath design is possible,
and you will use Type C tables for the rest of the course. Type E (exhaustive)
survives only while you have four inputs or fewer.

\subsection*{Equations}

\begin{align*}
\text{(1) canonical} \quad & Y_i = \overline{S_1}\,\overline{S_0}\,A_i + \overline{S_1}S_0\,B_i + S_1\overline{S_0}\,C_i + S_1S_0\cdot 0\\
& \qquad\qquad \text{row 3 outputs zero, so it contributes no product term}\\
\text{(2) factored} \quad & Y_i = \overline{S_1}\left( \overline{S_0}A_i + S_0 B_i \right) + S_1\overline{S_0}\,C_i
  \qquad \text{[distributive law, factoring $\overline{S_1}$]}\\
\text{(3) VHDL} \quad & \text{see \S4}
\end{align*}

\textbf{Look at what equation (2) is.} The bracket is a 2-to-1 multiplexer
choosing $A$ or $B$. The whole expression is a second 2-to-1 multiplexer choosing
between that result and $C$. You factored an equation and a \emph{circuit
structure} came out.

Factoring algebra and splitting a design into sub-blocks are the same operation.
This is the first place in the course you can see that, and it stays true all the
way to the CPU.

\section*{3. Guided examples}

\begin{description}[leftmargin=2.4em,style=nextline,itemsep=2pt]
\item[Example 1 --- three styles] The same 2-to-1 multiplexer written with
\code{when/else}, with \code{with/select}, and with \code{case} in a process.
One testbench asserts all three agree on every vector. Unlike Lab 2, the RTL
Viewer pictures are the same too --- here the difference is purely how the source
reads to a human.
\item[Example 2 --- \code{mux4}] A 4-to-1, 8 bits wide, generic. Four rows
specifying a function of 34 inputs.
\item[Example 3 --- the latch trap] Two 3-to-1 multiplexers, one of which leaves
\code{Sel = "11"} unassigned. Watch its output \emph{hold the previous value}
instead of taking a default, then find the latch in the RTL Viewer.
\end{description}

\textbf{Example 3 is the important one.} ``No assignment happens, so the output
keeps its previous value'' is Lab 1's latch, word for word. In Lab 1 that was the
feature. Here it is a bug: you wanted combinational logic and you described
something with memory in it.

\section*{4. Assignment}

\begin{enumerate}[leftmargin=2em,itemsep=3pt]
\item \textbf{\code{Mux2}} --- 2-to-1, generic width. Any of the three styles is
acceptable; say in your report \emph{why} you picked yours. That is the graded
part, not the syntax.

\item \textbf{\code{Mux3}} --- 3-to-1, generic width, \textbf{including row 3}.

Use a \code{case} with \code{when others}, not an \code{if/elsif} chain. VHDL
forces a \code{case} to cover every possible selector value, so the language
itself prevents the Example 3 mistake. An \code{if/elsif} chain will happily let
you forget.

\item \textbf{\code{CPUMuxes}} --- the four multiplexers the CPU needs, built by
instantiating your generic entities at the right widths.
\code{MUX2x5} is done for you as a worked example; write the other three the same
way.

\begin{center}
\begin{tabular}{llll}
\toprule
entity & width & sel & role on the CPU block diagram \\
\midrule
\code{MUX2x5}  & 5  & 1 & \code{DselMUX} --- register write-destination select \\
\code{MUX3x5}  & 5  & 2 & (spare 5-bit selector) \\
\code{MUX2x32} & 32 & 1 & \code{ALUinMUX} --- ALU A-input: register or immediate \\
\code{MUX3x32} & 32 & 2 & \code{DinMUX} --- write data: ALU / memory / return \\
\bottomrule
\end{tabular}
\end{center}

\item \textbf{Testbench} --- complete \code{tb\_Lab3.vhd}. \textbf{At least six
checked vectors}, including all of:
\begin{itemize}[leftmargin=1.4em,itemsep=1pt]
\item every selector value on \code{MUX3x32}, \textbf{including \code{"11"}};
\item a \textbf{no-latch check}: hold the selector, change the selected input,
and prove the output moves immediately. A latched output would keep the old
value --- and a latch is exactly what you get if you forget row 3;
\item at least one check on a 5-bit mux, proving the generic works at a width
other than 32.
\end{itemize}
\end{enumerate}

\subsection*{Coverage argument for a Type C table}

You cannot test $2^{98}$ combinations and you do not need to. The argument is:

\begin{quote}
\textbf{exhaustive on the selector, distinct patterns on the data.}
\end{quote}

``Distinct'' matters. If \code{A}, \code{B} and \code{C} all carry the same value,
a multiplexer that ignores its selector entirely and always returns \code{A}
passes every check. Give each input a value only it has --- \code{x"AAAAAAA1"},
\code{x"BBBBBBB2"}, \code{x"CCCCCCC3"} --- and a wrong selection cannot hide.

\section*{5. Evidence to collect (E1, E2, E7)}

\begin{center}
\begin{tabular}{lll}
\toprule
& Source & What it must show \\
\midrule
E1 & Flow Summary + Messages & 0 errors; \textbf{no inferred-latch warning} \\
E2 & RTL Viewer & multiplexers and \emph{no storage elements} \\
E7 & Questa waveform + transcript & every selector value, 0 assertion failures \\
\bottomrule
\end{tabular}
\end{center}

The top-level entity is \code{MUX3x32}. \textbf{Its register count must be zero.}
Any storage element at all means you inferred a latch --- go back to Example 3.

\section*{6. Report --- \code{lab03.tex}}

All seven sections. Specifically:

\begin{itemize}[leftmargin=1.6em,itemsep=1pt]
\item \textbf{Specification:} the Type C table, and one sentence on why it is
Type C and not Type E.
\item \textbf{Design:} the factored equation, the law named, and the observation
that factoring produced a multiplexer \emph{tree}. Also: which coding style you
chose and why.
\item \textbf{Results:} an \textbf{annotated copy of the CPU block diagram} with
your four multiplexers circled and labelled. Full page. This is the first time
you see the whole datapath, and from Lab 7 onward you will be reading it
constantly.
\item \textbf{Discussion:} required. One bug, how you found it.
\end{itemize}

\section*{7. Grading (100 points)}

\begin{center}
\begin{tabular}{lr}
\toprule
Pre-lab, submitted at the start of the session & 5 \\
Design compiles cleanly, warnings explained & 25 \\
Testbench and test vectors, coverage justified & 20 \\
Simulation results match the specification & 20 \\
Report chapter quality & 15 \\
In-lab demonstration and explanation & 15 \\
\midrule
Total & 100 \\
\bottomrule
\end{tabular}
\end{center}

Late work is accepted for one week at $-20$ points.

\section*{8. Things that will go wrong}

\begin{itemize}[leftmargin=1.6em,itemsep=1pt]
\item \textbf{Forgetting row 3.} You get a latch, a warning you did not read, and
a design that remembers things it should not.
\item \textbf{Missing \code{when others}.} \code{std\_logic} has nine values, not
two. The compiler will refuse the file --- which is the language doing you a
favour.
\item \textbf{\code{generic map} / \code{port map} order.} \code{generic map}
first, no semicolon between them.
\item \textbf{Identical test data on every input.} A mux that ignores its
selector passes. Use distinct patterns.
\item \textbf{Width mismatch in \code{port map}.} Usually a forgotten
\code{generic map}, reported at a confusing place.
\end{itemize}

\vspace{10pt}
\hr
{\small Every one of these four multiplexers is instantiated inside your CPU in
Lab 12. When you wire the datapath together, you will be connecting components
you wrote today.}

\end{document}
